\documentclass[11pt]{article}
\usepackage{series}

\renewcommand{\PartRoman}{VI}
\hypersetup{
  unicode=true,
  pdftitle={Fixed Points VI: Formal Synthesis and Physical Interpretations},
  pdfauthor={Peter Nero}
}

\title{Fixed Points VI: Formal Synthesis and Physical Interpretations}
\author{Peter Nero}
\date{September 2026\\Version 7}

\begin{document}
\maketitle
\seriestagline

\begin{abstract}
We synthesize the canonical FP--I--V results without promoting their control
parameter, projections, or diagnostics into an unproved fundamental field
theory. The rigorous spine consists of conditional projected fixed-point
existence, strict-Lyapunov promotion to equilibrium, joint-mode damping and
disturbance floors, perturbative persistence of a curved spectral cluster,
curvature leakage, intrinsic first-order centroid modulation, frozen linear
Ornstein--Uhlenbeck covariance, canonical-correlation bounds, and declared
admissibility exits. We distinguish three logical levels throughout:
inherited theorem, conditional model completion, and physical interpretation.
All inherited results are restated as self-contained contracts and cited to
their canonical paper; FP~VI does not reproduce their theorem/proof blocks.
Its own formal results are limited to the bilocal-causality obstruction and
the local-mediator completion.
A Lorentzian gauge/gravity action, quantum covariance, particle identities,
merger, measurement, and cosmology belong to the latter two levels unless
additional source and equivalence theorems are supplied. The Gaussian
Lyapunov sign and nonnormal resolvent bounds are stated in their domain-valid
forms, and we show why an
instantaneous equal-time bilocal kernel is not microcausal merely because it
leaves the principal symbol unchanged. A local mediator is the consistent
route to causal overlap dynamics. The synthesis uses only the preceding
Fixed Points papers and standard mathematical literature. Physical
interpretations are conditional applications, not additional proof inputs.
\end{abstract}

\section*{Revision note for version 7}
\begin{description}
\item[Supersedes.] Version~6, retained as the released edition.
\item[Reason.] Even an explicitly separate downstream progress ledger placed
later research inside a foundational synthesis and introduced reverse citations.
\item[Resolution.] The ledger and application-specific carrier claims are
removed. The source chain is now FP~I--V only, alongside standard mathematics;
the cited FP editions are aligned with the present foundational revisions.
\item[Retained result.] The analytic synthesis, its stated assumptions, and
the two bilocal-causality propositions retain their roles.
\item[Remaining boundary.] Concrete physical dynamics and their relation to
the control framework must be established by applications of this foundation.
\end{description}

\section*{Revision note for version 6}
\addcontentsline{toc}{section}{Revision note for version 6}
\begin{description}
\item[Supersedes.] \emph{Fixed Points VI: Formal Synthesis and Physical
Interpretations}, version~5.
\item[Reason.] Version~5 correctly separated the FP spine from later q79
work, but the synthesis still needed a clearer reading route and firmer
protection against treating restated contracts as duplicate theorem sources.
\item[Resolution.] Version~6 organizes the paper into rigorous, downstream,
and interpretive reading lanes; explains why contracts are repeated while
proofs are not; and preserves FP~I--V as the canonical owners of their
results. It also separates release description, revision history, and
computational provenance.
\item[Retained result.] The inherited FP~I--V spine and the two FP~VI
bilocal-causality results retain exactly their version~5 scope.
\item[Remaining boundary.] The selected action, physical q79 continuum
operator, general Born source, fixed-coupling QFT, global state, matching, and
cosmology remain downstream obligations.
\end{description}

\section*{Revision note for version 5}
\addcontentsline{toc}{section}{Revision note for version 5}
\begin{description}
\item[Supersedes.] \emph{Fixed Points VI: Formal Synthesis and Physical
Interpretations}, version~4.
\item[Reason.] Version~4 correctly narrowed the FP series, but it predates the
exact q79 lane-factor, projective-module, connection-compiler and nonlinear
strain-quotient results, as well as later scoped QM and QFT closures. Without a
separate status layer, ``not derived by FP'' could be misread as ``not achieved
anywhere downstream.''
\item[Resolution.] Version~5 preserves the FP theorem boundary, corrects the
q79 preprojection types, and adds an independently sourced downstream ledger
with exact closure counts and nonpromotion clauses. It also restores strict
theorem ownership: FP~I--V retain their respective fixed-point, damping,
curvature, and covariance theorems, while FP~VI owns only its two new
bilocal-causality results.
\item[Retained result.] The complete FP--I--V fixed-point, damping, curvature,
covariance, and admissibility spine survives at its declared theorem tier.
\item[Remaining boundary.] Selected HYM endpoints and action, the physical q79
Hessian and reduction, universal apparatus and actualization, a fixed-coupling
continuum QFT, selected global state, RG and matching, and cosmology remain
downstream programs.
\end{description}

\section*{How to use this synthesis}
\addcontentsline{toc}{section}{How to use this synthesis}
This paper is a map of the Fixed Points chain, not a substitute for its five
source papers.  It restates the input--output contracts needed to understand
the analytic sequence and cites each earlier canonical owner. Its only new
formal results concern the causality
failure of instantaneous bilocal overlap and the local-mediator completion.

\paragraph{The central picture in plain language.}
The series builds a sequence of conditional gates.  Geometry and spectral
data define a coherent projector.  Compactness or condensation gives a
projected return.  A strict Lyapunov law promotes that return to equilibrium.
Spectral margins control disturbances; curvature can rotate the projector and
source leakage; covariance describes a frozen linear fluctuation cloud; and
declared margins detect exit from a chosen admissible region.  Each gate has
its own hypotheses, and passing one does not silently pass the next.

\paragraph{Reading route.}
Sections~1--7 assemble the analytic results and their input contracts.
Sections~8--9 discuss an explicitly conditional local action and the
causality of overlap interactions. Section~10 separates possible physical
interpretations from the inherited mathematics, and Section~11 summarizes
the full dependency chain. No application-specific result is required to
read or establish that chain.

\paragraph{Why contracts are repeated but proofs are not.}
A standalone synthesis must state enough hypotheses and conclusions to be
understood without opening five documents at once.  It therefore repeats
short contracts, not theorem ownership.  Proofs and canonical formal
statements remain in FP~I--V.  This prevents the same theorem from acquiring
multiple apparent sources while preserving standalone readability.

\paragraph{Scope boundary.}
The control flow, classical covariance, and admissibility diagnostics do not
by themselves define Lorentzian dynamics, quantum commutation relations,
particle identities, or cosmology. Applications may use these results after
checking their hypotheses; their additional conclusions are not premises
of the Fixed Points theorems.

\tableofcontents

\section{Logical levels and common geometry}

Every statement in this synthesis has one of the following statuses.
\begin{description}
\item[Inherited theorem.] A consequence of FP~I--V under hypotheses restated
or cited here.
\item[Conditional completion.] A mathematically specified extra model whose
conclusions hold if that model and its hypotheses are adopted.
\item[Physical interpretation.] A proposed reading that is not established by
the fixed-point theorems alone.
\end{description}

Let $\mathcal H$ be the declared control Hilbert space. The joint internal
operator is constructed from three declared vertical operators using the
strong-commutation and domain hypotheses of FP~II. Circle--lens--nil labels
refer here to obstruction, rank, or operator data; they do not assert a literal
product or nesting of three manifolds. A common $U(1)$ phase may be represented
by a line connection without adding a product dimension, as in FP~II.
Finite rank flags do not specify differential operators or a Galerkin basis.
Any such realization must be supplied before invoking the analytic estimates;
no selected compactification or physical carrier is assumed here.

The coherent projector $P$ is the joint spectral projector. A product of three
bundle projectors is permitted only when those projectors strongly commute;
otherwise the product need not be an orthogonal projector. Write $Q=I-P$.

On a curved bundle, the normalized Laplace-type operator is
\[
 L_x=\nabla_x^\ast\nabla_x+\mathcal R_x.
\]
When the selected low cluster remains separated, FP~IV defines the curved
Riesz projector $P_{\mathcal R}(x)$. Retaining the old projector $P$ instead
requires the off-diagonal leakage block $Q\mathcal RP$.

The parameter of the FP gradient flow is a stabilization or control
parameter. Identifying it with Lorentzian proper time, coordinate time, or a
quantum evolution parameter is a conditional completion, not an inherited
theorem.

\section{Fixed-point and equilibrium synthesis}

Let $\Phi_\tau$ be a well-defined time-$\tau$ map for the declared control
flow and set $T=P\Phi_\tau$ on a closed bounded convex set $K\subset\mathcal H$.

\paragraph{Imported fixed-point chain.}
Fixed Points~I and II are the canonical sources for this chain
\cite{FPI,FPII}. For standalone use, its complete contract is as follows.
Assume one of the following FP~I alternatives:
\begin{enumerate}
\item $T(K)\subset K$ and $T:K\to K$ is continuous and compact; or
\item $T(K)\subset K$ and $T$ is continuous and condensing for a declared
measure of noncompactness.
\end{enumerate}
Then $T$ has a fixed point $u_\ast\in K$. This proves
$P\Phi_\tau(u_\ast)=u_\ast$. If, in addition, the trajectory through
$u_\ast$ remains in the projected invariant set and obeys a strict Lyapunov
identity
\[
 \mathcal C(u(\tau))-\mathcal C(u(0))
 =-\int_0^\tau\mathcal D(u(s))\,ds,
 \qquad \mathcal D\ge0,
\]
with $\mathcal D(u)=0$ exactly at equilibria, then $u_\ast$ is an equilibrium
of the declared flow.
The first conclusion is an application of Schauder's or Darbo--Sadovskii's
theorem. At a
projected recurrent point, the Lyapunov values at the two endpoints agree.
The displayed identity forces $\mathcal D=0$ along the intervening orbit, and
strictness gives equilibrium. Without the second step, a projected time-step
fixed point is not automatically a steady solution.

\paragraph{Imported contraction-uniqueness gate.}
Fixed Points~I owns the generic Banach gate. In the steady-equation form used
here, suppose a steady equation can be written on a specified Banach space as
$u=-A^{-1}N(u)$, where $A^{-1}$ exists on that space and
\[
 \|A^{-1}\|\operatorname{Lip}(N)<1
\]
on an invariant complete subset. Then the steady solution there is unique.
No uniqueness conclusion follows from a vertical spectral gap if $A$ has an
unremoved kernel or if the displayed inverse estimate is not proved.

\section{Joint-mode stability and disturbance floors}

Let $\alpha$ be one joint modal label, including multiplicity. After the
FP~III remainder estimate \cite{FPIII}, suppose a noncoherent amplitude satisfies
\[
 \frac{d}{dt}|a_\alpha|
 \le-\gamma_\alpha|a_\alpha|+|f_\alpha(t)|,
 \qquad \gamma_\alpha>0.
\]
Then the deterministic input-to-state floor is
$\|f_\alpha\|_\infty/\gamma_\alpha$. For the exact scalar It\^o equation
\[
 da_\alpha=-\gamma_\alpha a_\alpha\,dt
 +\sqrt{q_\alpha}\,dW_\alpha,
\]
the stationary variance is $q_\alpha/(2\gamma_\alpha)$. Deterministic force
amplitude and stochastic power have different units and cannot be represented
by one shared disturbance threshold.

For infinitely many modes, bundlewise stochastic control requires the
appropriate weighted covariance trace; deterministic control requires its
own weighted series. These are sufficient nonlinear stability criteria.
Gaussian invariant laws are asserted only for exact linear OU systems.

\section{Curvature, centroid motion, and interaction}

Use the curved-operator framework of FP~IV \cite{FPIV}. Let the uncurved
noncoherent spectrum start at $\lambda_\ast>0$. If
$\|\mathcal R\|<\lambda_\ast/2$, the selected low spectral cluster remains
separated and has a curved Riesz projector of the same rank. If the uncurved
projector is retained and the one-sided $Q$-sector damping margin is
$\gamma_Q>0$, FP~IV gives the schematic sharp form
\[
 \frac{d}{dt}\|q\|
 \le-\gamma_Q\|q\|+\|Q\mathcal RP\|\,\|p\|+\|QF\|,
\]
and therefore a curvature-induced leakage floor.

For a localized profile inside a strongly convex normal ball, the centroid is
the Karcher mean. A first-order gradient parent flow yields a first-order
modulation equation
\[
 G_{ij}(X)\dot X^j=-\partial_iV_{\rm eff}(X)+\varepsilon_i.
\]
A Newton equation requires a separately declared inertial parent theory.

For two profiles, absolute cross-term control proves only
\[
 |E_{\rm int}|\le C\mathcal O.
\]
Attraction requires $E_{\rm int}\le-c\mathcal O$; repulsion requires the
opposite sign. Even a proved attractive sign does not establish merger,
collapse, or a transition rate without a dynamical connection theorem.

\section{Frozen Gaussian multi-structure theory}

Consider the real linear SDE
\begin{equation}
 d\zeta=A\zeta\,dt+B\,dW_t,
 \qquad Q:=BB^\top .
 \label{eq:linear-sde}
\end{equation}

\paragraph{Imported semigroup-covariance result.}
Fixed Points~V is the canonical source for the nonnormal covariance theorem
\cite{FPV}. For the real linear SDE~\eqref{eq:linear-sde}, assume
\[
 \|e^{tA}\|\le M e^{-\omega t},
 \qquad t\ge0,\quad M\ge1,\quad\omega>0.
\]
Then~\eqref{eq:linear-sde} has stationary covariance
\[
 \Sigma=\int_0^\infty e^{tA}Qe^{tA^\top}\,dt,
\]
which solves
\[
 A\Sigma+\Sigma A^\top+Q=0.
\]
Moreover,
\[
 \|\Sigma\|\le\frac{M^2\|Q\|}{2\omega},
 \qquad
 \|A^{-1}\|\le\frac{M}{\omega}.
\]
The semigroup estimate makes the covariance and resolvent integrals converge.
Differentiation of the
covariance integrand gives the Lyapunov equation, while
$A^{-1}=-\int_0^\infty e^{tA}\,dt$. Taking norms proves the bounds.

For a normal $A$, one may often take $M=1$ with $\omega$ given by the spectral
abscissa. For a nonnormal matrix, the spectral abscissa alone does not imply
$\|A^{-1}\|\le1/\omega$ or
$\|\Sigma\|\le\|Q\|/(2\omega)$; transient amplification must be controlled.

Under the block-diagonal damping assumptions of FP~V, the cross covariance
satisfies a Sylvester equation. Canonical correlation is normalized by the
smallest positive eigenvalues of the marginal covariance blocks, not their
largest operator norms.

\section{Classical covariance versus quantum covariance}

The word ``covariance'' occurs in both classical probability and quantum
theory, but the objects satisfy different constraints.  The FP~V Lyapunov
equation produces ordinary second moments of a classical stochastic process.
A quantum covariance becomes physical only after a symplectic or canonical
commutation structure and compatible completely positive dynamics have been
specified.

The covariance produced by~\eqref{eq:linear-sde} is classical. To interpret it
as a bosonic quantum covariance one must separately specify canonical
commutation relations, a value of $\hbar$, and quantum dynamics whose noise
and damping satisfy the complete-positivity constraints. Only then is
\[
 \Sigma+\frac{i\hbar}{2}\Omega\succeq0
\]
the physical uncertainty condition. Classical positivity of $\Sigma$ does not
imply this inequality.

For a declared bipartite Gaussian quantum state, nonpositivity after partial
transpose certifies entanglement. PPT is exact for the two-mode $1\times1$
case \cite{SimonPPT}; in general multimode bipartitions PPT need not be sufficient for
separability. These facts classify an already constructed quantum Gaussian
state. They do not quantize the FP control flow or derive particle statistics.

\section{Admissibility and exit}

Let $m_j(x)$, $j\in\mathcal J$, be a finite family of continuous declared
margins: curved cluster separation, one-sided damping, or explicit
deterministic/stochastic performance tolerances. Define
\[
 \mathcal D_\varepsilon
 =\{x:m_j(x)\ge\varepsilon\ \forall j\},
 \qquad
 \mathfrak B(x)=\max_j[-m_j(x)].
\]
Then $x\in\mathcal D_0$ exactly when $\mathfrak B(x)\le0$. Along a continuous
path, a change from $\mathfrak B\le0$ to $\mathfrak B>0$ has a first boundary
time at which $\mathfrak B=0$.

This is an exact exit diagnostic for the declared constraints. It does not
prove that $\mathfrak B$ is a physical energy or force. It does not establish
which state follows exit. A mountain-pass theorem needs an actual energy
functional, and basin selection needs a post-exit evolution or transition
kernel.

For affine observables of a frozen Gaussian system, FP~V supplies finite-grid
and Borell--TIS continuous-time exit estimates. A Lipschitz nonlinear function
of a Gaussian process is not generally Gaussian. Equal-time cross-correlation
can bound simultaneous observed exits, but cannot prove causal propagation or
non-propagation.

\section{Status of a Lorentzian master action}

This section illustrates a conditional completion.  Writing a familiar local
action shows what additional fields and equations could supply physical
evolution; it does not show that the FP control data select that action.  The
gap between a writable action and a derived action is precisely the missing
source and matching problem.

A local Lorentzian gauge/gravity/matter action may be appended as a
\emph{conditional completion}, for example schematically
\[
\begin{aligned}
S_{\rm cand}=\int\sqrt{-g}\,\Bigl[
&\frac{M_{\rm Pl}^2}{2}R-\Lambda
-\frac14\sum_r\operatorname{tr}F_r^2\\
&-\sum_a|D\varphi_a|^2-V(\varphi)
+\sum_b\bar\psi_b(i\gamma^\mu D_\mu-M_b)\psi_b
\Bigr]\,d^4x.
\end{aligned}
\]
The FP~I--V theorems do not select the gauge group, representations, number of
fields, couplings, Yukawa matrices, or potential. Nor do they prove that the
gradient control flow is the Euler--Lagrange evolution of $S_{\rm cand}$.

If a mass term depends on curvature, such as
$M_a^2=m_a^2+\xi_aR$, its metric variation contributes nonminimal terms to
the gravitational field equation; it cannot be inserted while retaining the
minimally coupled Einstein equation unchanged. Likewise, a point-particle
Lorentz-force law requires a controlled localized-solution limit and is not a
consequence of gauge covariance alone.

Three internal structures do not by themselves prove three gauge factors,
bosonic or fermionic statistics, bifundamental matter, the Standard Model
particle assignment, anomaly cancellation, or equivalence with a quantum
field theory. Those require independent representation, source, and matching
theorems.

\section{Bilocal overlap and causality}

Consider an equal-parameter interaction of the form
\[
 \mathcal N[\phi](t,x)
 =\int_{\Sigma_t}K(x,y)|\phi(t,y)|^2\,d\mu_t(y)\,\phi(t,x).
\]

\begin{proposition}[Instantaneous bilocal obstruction]\label{prop:bilocal}
If $K(x,y)$ is nonzero for spatially separated $x$ and $y$, the equation at
$(t,x)$ depends instantaneously on data at $y$. Leaving the differential
principal symbol unchanged is therefore insufficient to prove the usual
local domain-of-dependence property. An $L^1$ bound on $K$ may help local
well-posedness, but it does not establish microcausality.
\end{proposition}

\begin{proof}
Choose two initial data sets that agree near $x$ but differ near a separated
point $y$ where $K(x,y)\ne0$. Their bilocal source values at $(t,x)$ differ at
the same parameter time. Thus the source is not determined by data in an
arbitrarily small causal neighborhood of $x$.
\end{proof}

\begin{proposition}[Local-mediator completion]\label{prop:mediator}
Introduce a local mediator $\chi$ with a hyperbolic equation, for example
\[
 (\Box_g+m_\chi^2)\chi=g_\chi|\phi|^2,
\]
and couple $\chi$ locally back to $\phi$. Under the standard regularity,
gauge, and constraint hypotheses for the resulting local hyperbolic system,
its domain of dependence is governed by the local characteristics. Eliminating
$\chi$ with a retarded Green operator produces a causal retarded memory
kernel, not in general a symmetric equal-time bilocal action.
\end{proposition}

This local parent theory is the preferred route if overlap dynamics is meant
to represent physical causal interaction.

\section{Physical interpretation ledger}

The ledger should be read as a type system for claims.  ``Supported'' means
the stated control model and hypotheses already imply the claim.
``Conditional'' means an additional mathematical model would make the claim
well posed and testable.  ``Not derived'' means the Fixed Points chain alone
contains no implication to that conclusion; it is not a declaration that the
idea is impossible or absent from every other MTT paper.

\paragraph{Supported within the declared control model.}
The declared model supports projection into coherent and noncoherent sectors
under operator-domain hypotheses. It also supports conditional fixed points,
strict-Lyapunov equilibrium promotion, damping estimates and disturbance
floors, perturbative curved-cluster persistence, curvature leakage, intrinsic
first-order centroid modulation, frozen linear covariance, and admissibility
exit diagnostics.

\paragraph{Conditional after adding a model.}
Gauge forces, Lorentzian motion, quantum uncertainty, Gaussian entanglement,
decoherence, thermodynamic entropy production, and causal overlap can be
studied after a compatible local Lorentzian or open-quantum model is specified.
Their parameters and consistency conditions are additional data.

\paragraph{Not derived by this series.}
The present theorems do not derive boson/fermion identity, particle masses,
Standard Model representations, merger or collapse rules, the Born rule,
measurement outcomes, emergent time, inflation, horizons, or a cosmological
arrow of time. This is a statement about the implications of this series,
not a claim about the status of every physical theory or application.

\section{Synthesis ledger and remaining obligations}

\paragraph{Imported FP synthesis.}\mbox{}\\
Assume the operator, compactness or condensing, strict-Lyapunov, one-sided
damping, curvature-smallness, covariance, and margin hypotheses stated in
FP~I--V and summarized above. Then the declared control model has:
\begin{enumerate}
\item a projected time-step fixed point, promoted to an equilibrium only by
the strict Lyapunov condition;
\item quantified noncoherent damping and deterministic/stochastic floors;
\item a persistent curved low cluster and a quantified old-projector leakage
term;
\item intrinsic first-order centroid modulation under the modulation
hypotheses;
\item exact stationary covariance for the frozen linear OU sector and
correctly normalized correlation bounds; and
\item exact detection, plus conditional Gaussian probability bounds, for
exit from a declared admissible domain.
\end{enumerate}
No Lorentzian, quantum, particle-physics, merger, measurement, or cosmological
claim follows without its additional completion theorem.

For an application, the constructive obligations are to specify an energy and
its operator domain, establish the invariant set and relevant analytic bounds,
and prove the relation between the control flow and the proposed physical
evolution. Quantum or relativistic interpretations need their own state,
causality and observable structures. The present paper neither supplies those
structures nor needs them to establish the analytic synthesis.

\section{Conclusion}

The corrected Fixed Points series is a coherent conditional control and
spectral framework. Its strongest results concern existence under explicit
compactness hypotheses, stability margins, curved spectral persistence,
leakage, modulation, Gaussian covariance, and admissibility diagnostics. FP~VI
does not weaken those results by surrounding them with broader interpretations;
it identifies exactly what they prove and exactly which bridges remain before
they can support a fundamental spacetime, quantum, or particle theory.

The synthesis also clarifies how to use the series.  FP~I owns the analytic
fixed-point foundation; FP~II owns the ten-dimensional joint-operator
specialization; FP~III owns disturbance floors and conditional
homogenization; FP~IV owns curvature, finite reduction, centroid, and
transition boundaries; FP~V owns multi-structure covariance and
admissibility-exit probabilities. FP~VI connects those contracts and adds
only the bilocal-causality analysis.
The next advances must therefore instantiate the missing source data or
physical completion, not restate the existing fixed-point chain under a new
heading.

\begin{thebibliography}{99}

\bibitem{FPI}
P.~Nero,
\newblock \emph{Fixed Points I: Fixed Points over Multi--Bundle Manifolds},
\newblock version~7, 2026.

\bibitem{FPII}
P.~Nero,
\newblock \emph{Fixed Points II: Projected Fixed Points and Equilibria in a
10D Modal Model},
\newblock version~6, 2026.

\bibitem{FPIII}
P.~Nero,
\newblock \emph{Fixed Points III: Disturbance--Damping Balance and Stability},
\newblock version~8, 2026.

\bibitem{FPIV}
P.~Nero,
\newblock \emph{Fixed Points IV: Curvature, Centroid Motion, and Structural
Transitions on Bundle Manifolds},
\newblock version~7, 2026.

\bibitem{FPV}
P.~Nero,
\newblock \emph{Fixed Points V: Curvature Coupling, Multi--Structure Dynamics,
and Admissibility Barriers},
\newblock version~9, 2026.

\bibitem{DaPratoZabczyk}
G.~Da~Prato and J.~Zabczyk,
\newblock \emph{Stochastic Equations in Infinite Dimensions},
\newblock Cambridge University Press, 1992.

\bibitem{SimonPPT}
R.~Simon,
\newblock Peres--Horodecki separability criterion for continuous variable
systems,
\newblock \emph{Physical Review Letters} \textbf{84} (2000), 2726--2729.

\end{thebibliography}

\section*{Dependency and reproducibility statement}
The preceding FP papers are the sole MTT proof sources used here. Standard
mathematics and the two local causality arguments complete the synthesis.
No later research ledger, numerical packet or physical endpoint is needed.

\end{document}
